%% ============================================================
\section{Introduction}
\label{sec:intro}
%% ============================================================

%% --- Domain landscape ---

Large language models are increasingly deployed not as isolated assistants but as autonomous agents that conduct entire research workflows---from literature review and hypothesis generation through experimental design and execution to manuscript writing~\citep{lu2024aiscientist, yamada2025aiscientistv2, huang2025deepscientist, tang2025airesearcher, schmidgall2025agentlab, jansen2025codescientist}.
On systems-optimization tasks, such agents now produce solutions competitive with human experts~\citep{cheng2025barbarians, novikov2025alphaevolve}, and end-to-end pipelines have generated papers accepted at peer-reviewed workshops~\citep{yamada2025aiscientistv2}.
The resulting artifacts---code, experimental results, and professional-looking manuscripts---are increasingly difficult to distinguish from human-authored research on surface quality alone.

%% --- Problem buildup: the verifiability gap ---

This rapid capability growth exposes a structural tension between \emph{generation} and \emph{verification}.
Autonomous research systems operate as multi-stage pipelines in which each stage consumes the output of the previous one: a literature summary shapes the hypothesis, the hypothesis determines the experiment, and experimental results feed into the manuscript.
In such architectures, errors introduced at any stage are not merely preserved but amplified---a flawed summary can bias experimental design, and a misinterpreted result can carry through into a paper that appears internally coherent, precisely because the same error is reflected consistently across sections.
The risk grows with trajectory length: agents struggle to track an ever-expanding context \citep{liu2024lost,liu2023agentbench}, hallucinate, and drift from the original objective.
The problem is exacerbated by fundamental limitations in how language models handle evidence: generated text is difficult to verify against sources~\citep{liu2023evaluating}, factual claims drift from their grounding~\citep{min2023factscore}, and scientific citations are frequently inaccurate or fabricated~\citep{press2024citeme}.

In autonomous pipelines, these failure modes interact and compound---a model can overstate method descriptions beyond what the code implements, report scores that do not reproduce under the benchmark's own evaluator, and populate bibliographies from parametric memory rather than retrieval, all while producing text that reads as technically sound.
Existing evaluation protocols, whether automated review scores or benchmark leaderboards, assess surface presentation (i.e., how the paper reads) and procedural completion but do not check whether individual claims trace to supporting evidence.

%% --- Condensed motivation from our audit ---

This verifiability gap is not hypothetical.
In a systematic audit of 75~papers from five autonomous research systems across five benchmark tasks, we find that \emph{every baseline system exhibits evidence chain failures}: hallucinated references that do not correspond to any real publication (up to 21\% of all bibliography entries), method sections that describe algorithms not present in the submitted code, unreproducible scores, and solution code that exploits the evaluator rather than solving the task.
These failures share a common root cause: \emph{no existing evaluation protocol audits whether claims are supported, and no existing autonomous research system is designed to trace claims back to evidence.}

%% --- Our answer: CoE ---

We address this with \emph{\coe{}} (CoE), a verifiability framework for AI-driven research.
Just as ACID\footnote{Atomicity, consistency, isolation, durability.}~\citep{haerder1983principles} defines what ``reliable'' means for a database transaction,
CoE defines what ``verifiable'' means for a research claim: \textbf{every claim must trace, through a recorded evidence chain, to a grounding source.}
We instantiate CoE in three ways:

\begin{enumerate}[nosep,leftmargin=*]
  \item \textbf{The CoE Standard} (\S\ref{sec:coe}): a claim taxonomy (citation, numerical, methodological, conclusion) and the evidence chain structure required for each type.

  \item \textbf{\sys{}} (\S\ref{sec:system}): an end-to-end autonomous research system whose pipeline---Problem Investigator, Discovery Engine, and Paper Writer with Claim Verifier---is designed to satisfy CoE natively.
  The Problem Investigator reads up to 100 full-text PDFs per topic, producing grounded experiment briefs. And the Claim Verifier checks every claim in the draft against its declared evidence source before the final paper is produced.

  \item \textbf{\coeaudit{}} (\S\ref{sec:coeaudit}): a post-hoc audit for verifying an AI-driven research paper through four integrity checks---Score Verification, Specification Violation, Reference Verification, and Method-Code Alignment---targeting the most damaging evidence chain failures.
\end{enumerate}

We apply \coeaudit{} to 15~papers from each of five systems across five frontier systems-research tasks from ADRS~\citep{cheng2025barbarians,skydiscover2026}~(\S\ref{sec:experiments}).
Every baseline exhibits at least one integrity check failure.
\sys{} achieves zero hallucinated references (0/337 bibliography entries), perfect score verification (12/12), and the highest method--code alignment (14/15), while matching or exceeding human expert solver performance on all five tasks.
We further demonstrate that \sys{} generalizes to six additional tasks spanning medical imaging, fine-grained recognition, 3D perception, and parameter-constrained language modeling, achieving state-of-the-art on Parameter Golf and gold medals on MLE-Bench tasks where baselines fail entirely.